\chapter{The JML-Inferrer System}
\label{ch:inference}

This chapter presents JML-Inferrer, the static-analysis system that produces
every specification used in the remainder of the thesis. It is the
load-bearing artefact of the work: the empirical studies of
Chapters~\ref{ch:testgen}, \ref{ch:distribution}, and~\ref{ch:compositional}
all consume its output, and the central methodological claim of the thesis ---
that heuristic inference validated against a formal oracle yields useful
specifications --- is a claim about this system. The chapter describes the
system at the level of detail required to interpret the empirical results that
follow: its architecture (Section~\ref{sec:inf-arch}), the
weakest-precondition/strongest-postcondition design discipline that organises
its inference (Section~\ref{sec:inf-discipline}), the inference of each clause
kind (Sections~\ref{sec:inf-pre} through~\ref{sec:inf-loops}), the
interprocedural propagation that handles method calls
(Section~\ref{sec:inf-inter}), the formal validation layer that recovers
soundness (Section~\ref{sec:inf-validation}), and the method categorisation
scheme used to structure the evaluation (Section~\ref{sec:inf-cat}).

\section{Architecture}
\label{sec:inf-arch}

JML-Inferrer is implemented in Java~21 and comprises approximately twenty
thousand lines of code at the time of writing. It uses
JavaParser~\cite{javaparser} to construct abstract syntax trees from Java
source and processes source trees recursively, inserting JML annotations in
place. The pipeline has three phases: parsing, specification inference, and
formatting. Parsing and formatting are standard --- the first builds an AST
from each compilation unit, the third serialises inferred specifications as
JML comments and writes the annotated source back --- and the analytical
substance is entirely in the second phase, which this chapter describes.
Average processing time is 88 milliseconds per file, a figure that places the
system comfortably within the budget of a continuous-integration step, a
pre-commit hook, or an interactive editor integration.

The inference phase is built on the visitor pattern over AST nodes. A
coordinating visitor traverses each compilation unit and, for every method
declaration that is eligible for processing, delegates to a specification
inferrer that analyses the method body and returns a structured specification
object holding the inferred preconditions, postconditions, loop invariants,
and frame conditions. The coordinating visitor then formats this object as JML
and injects it as a structured comment attached to the method. The design is
deliberately non-destructive: a method that already carries a specification ---
a hand-written \texttt{requires}, \texttt{ensures}, or
\texttt{loop\_invariant} --- is skipped, so the engine never overwrites a
human-supplied contract. The non-destructive principle reflects the engine's
intended role as an assistant rather than an authority: where a human has stated a
contract, that contract is taken as ground truth and left untouched, and the engine
fills only the gaps the human left. The principle also makes the engine safe to run
repeatedly over an evolving codebase --- a second run does not disturb the
specifications the first produced unless the code changed --- which is what allows
the continuous-integration deployment the thesis envisions, where the engine runs
on every change and refreshes only the specifications whose methods changed. The
formatting phase preserves any existing Javadoc, appending the JML rather than
replacing the documentation, and renders the inferred clauses legibly --- on
separate lines, simplified, in the conventional order of preconditions before
postconditions before frame conditions --- because a specification a reviewer
cannot read is one a reviewer cannot check, and the review queue for flagged
clauses depends on the emitted form being legible. Methods with no body (abstract
methods, interface methods) and methods whose specifications are inferred elsewhere
are handled specially, a point that becomes important for the distribution work of
Chapter~\ref{ch:distribution}, where interface methods turn out to be a source
of subtle bugs.

A design constraint maintained throughout the implementation is a file-size
ceiling: no source file exceeds five hundred lines, and a class that grows
beyond it is split into focused sub-analysers by logical grouping ---
precondition analysis, postcondition analysis, loop-invariant analysis, frame
analysis, and so on, composed by delegation rather than gathered into a
monolith. The constraint is a software-engineering discipline rather than an
analytical one, but it matters for the maintainability of an engine whose
heuristic repertoire grows over time, and it reflects the broader stance of
the project that the inferrer is itself a piece of engineered software held to
the standards it helps enforce on others.

\subsection{Design Philosophy and Engineering Constraints}

The engine is a research instrument, and its construction was governed by a
philosophy worth stating because it shaped decisions throughout. The governing
principle is that the correct solution is preferred to the easy one: when the
engine produces invalid JML, the response is to fix the inference to produce valid
JML, not to suppress the offending output. A pattern that cannot be expressed
soundly in JML is replaced by the closest valid expression that captures part of
its intent, and inference for a pattern is removed only as a last resort, when no
valid JML can be produced for it at all. This principle is what keeps the engine's
output a set of genuine, checkable specifications rather than a mixture of valid
JML and natural-language approximations, and it is enforced by the requirement that
every emitted clause be syntactically valid JML with no undefined variables and no
expression that is false at the point it is asserted.

The engineering constraints follow from the engine's status as evolving research
software. The five-hundred-line file ceiling, mentioned above, forces the
decomposition into focused sub-analysers. Every change to the inference logic is
accompanied by verification tests that prove the inferred specifications remain
valid, and a regression suite gates changes against the formal oracle, so that an
improvement to one pattern cannot silently break another. The testing requirement
is two-tiered, as Section~\ref{sec:inf-validation} describes, and the preference is
always for testing through the full inference pipeline rather than against
hand-written JML, because the engine's purpose is to \emph{infer} specifications
and a test that bypassed inference would not exercise that purpose. These
constraints are not incidental; they are what allow a heuristic engine, whose
pattern repertoire grows continually, to remain correct as it grows, and they are
the software-engineering counterpart to the methodological discipline of validating
every clause.

\subsection{Why Pattern Matching Rather Than Symbolic Execution}

The most consequential architectural decision is the deliberate avoidance of
symbolic execution~\cite{king1976symbolic, cadar2008klee} and full
weakest-precondition computation. Both require an SMT decision procedure and a
complete operational semantics for Java, and neither scales to arbitrary code:
symbolic execution suffers path explosion on branch-heavy code and diverges on
loops without a supplied bound, and a complete WP transformer needs the loop
invariants that are themselves the hardest thing to infer. The engine instead
performs pattern matching over AST nodes, recognising syntactic shapes that
correspond to specification clauses and emitting the corresponding JML. This
trades formal completeness for scalability and broad applicability --- an
established trade-off in industrial static analysis, where Infer and similar
tools accept unsoundness in exchange for the ability to run on millions of
lines~\cite{calcagno2015infer}.

The trade-off would be indefensible if it meant emitting unsound
specifications with no recourse. The recourse is the validation layer of
Section~\ref{sec:inf-validation}: the engine matches patterns aggressively and
discharges the results conservatively, so that the unsoundness inherent in
pattern matching is caught after the fact by a verifier rather than shipped to
the consumer. This two-layer architecture --- aggressive heuristic inference,
conservative formal validation --- is the system's answer to the soundness
objection, and it recurs as a theme throughout the thesis.

The trade-off has a precise cost and a precise benefit that are worth naming. The
cost is recall: a pattern matcher recognises only the shapes it has been taught,
so a precondition expressed by a guard idiom the engine does not recognise, or a
postcondition that depends on a computation the patterns do not reach, is missed.
The recall figures of Chapter~\ref{ch:testgen} measure this cost directly. The
benefit is scale and applicability: the engine runs in milliseconds per file on
any compilable Java, with no test suite, no documentation, and no SMT solver in the
inference loop, which is what makes it practical to run over twenty-thousand-method
libraries and to integrate into a build. A symbolic executor or a complete WP
transformer would recover some of the missed recall but at a cost in scale that
would foreclose the corpus-wide measurements the thesis depends on. The pattern-
matching choice is thus not a compromise forced by implementation expedience but a
deliberate position on the recall-versus-scale frontier, chosen because the thesis's
questions --- does inference help, can it be distributed, what does a second pass
add --- are questions about scale, answerable only by a tool that runs at scale,
and because the recall the choice sacrifices is partly recoverable by the
complementary techniques (dynamic detection, documentation extraction) the thesis
identifies as future ensemble partners.

\subsection{The Sub-Analyser Composition}

The five-hundred-line file ceiling forces a composition the inference engine
would benefit from in any case: rather than a monolithic inferrer, the engine is
a collection of focused sub-analysers, each responsible for one clause kind or
one pattern family, composed by delegation behind the coordinating visitor. A
precondition analyser owns the guard-and-throw recognition; a postcondition
analyser owns return and assignment structure; a loop-invariant analyser owns the
four loop heuristics; a frame analyser owns purity and assignable inference; an
overflow analyser owns the return-type gating of arithmetic clauses; an
interprocedural analyser owns the call-site propagation; a class-invariant
analyser owns field-level invariants; and a string analyser owns the
value-equality subtleties of \texttt{String}-returning methods. The coordinating
visitor invokes each in turn on a method and assembles their outputs into the
method's specification object.

The decomposition has a methodological payoff beyond maintainability. Because each
sub-analyser is independently testable, the engine's regression suite can target a
specific analyser's behaviour --- asserting, for instance, that the precondition
analyser emits exactly the expected \texttt{requires} clauses for a given guard
shape --- without running the whole pipeline. The suite is two-tiered: analysis
tests call a sub-analyser directly and assert on its inferred clauses without
invoking the verifier, while verification tests run the full pipeline through
OpenJML and assert that the inferred specification discharges. The two tiers
correspond to the two oracles of Section~\ref{sec:inf-validation} --- the analysis
tier checks what the engine emits, the verification tier checks that what it emits
is sound --- and the decomposition into sub-analysers is what makes the analysis
tier precise enough to localise a regression to a single pattern family. A change
that improves one analyser can be verified not to regress another, which matters
for an engine whose heuristic repertoire grows continually as new patterns are
added.

The deduplication of clauses is handled at assembly: each clause-adding operation
on the specification object checks for an existing equal clause before adding, so
that two sub-analysers that independently infer the same clause --- the
precondition analyser and the interprocedural analyser both concluding a parameter
is non-null, say --- do not produce a duplicate. This matters for the distribution
work of Chapter~\ref{ch:distribution}, where duplicate clauses would inflate the
embedded size, and for the readability of the emitted JML, where a duplicated
\texttt{requires} clause would be noise.

\section{The WP/SP Design Discipline}
\label{sec:inf-discipline}

The engine performs no symbolic WP or SP computation, but the WP and SP
calculi are not thereby irrelevant to it. They serve as a \emph{design
discipline}: they determine which AST patterns the engine recognises, what JML
clause each pattern produces, and how clauses propagate across procedure
boundaries. The discipline is what keeps the heuristic repertoire principled
rather than ad hoc --- every pattern the engine matches is justified by a
correspondence to a step of WP or SP reasoning, and patterns that have no such
correspondence are deliberately excluded.

The correspondence runs as follows. Preconditions are WP-derived: a
guard-and-throw idiom corresponds to the weakest precondition of the
normal-return path, and the engine emits the negated guard as a
\texttt{requires} clause. Postconditions are SP-derived: a return statement and
a field assignment each correspond to a strongest-postcondition consequence,
and the engine emits the corresponding equality. Interprocedural propagation
follows the textbook WP rule for procedure call: the callee's contract is
instantiated by formal-to-actual substitution at the call site. Loop invariants
are the point at which syntactic WP/SP reasoning is incomplete, and the engine
supplies heuristics in their place. The remainder of this chapter develops each
of these correspondences in turn, and in doing so makes precise what the engine
looks for and why.

\section{Precondition Inference}
\label{sec:inf-pre}

Preconditions arise, under the WP discipline, from guards that abort the normal
execution path. Consider a method body that begins with an early statement of
the form \texttt{if (G) throw \dots} or \texttt{if (G) return \dots}. Taking
the postcondition $Q$ to be the post of the normal-return path, a weakest-
precondition analysis propagates $\neg G$ backward to the method entry: the
precondition is whatever must hold on entry to avoid the abort. The engine
recognises throw-style guards, early-return guards, and validation-style
guards (calls into validation utilities such as \texttt{Objects.requireNonNull}
or a library's \texttt{Validate} class), and emits the negation of each guard
as a \texttt{requires} clause.

The most common shapes are worth enumerating because they account for the bulk
of inferred preconditions. A null guard \texttt{if (p == null) throw \dots}
yields the precondition \texttt{p != null}. A range guard
\texttt{if (i < 0 || i >= n) throw \dots} yields two conjuncts, \texttt{i >= 0}
and \texttt{i < n}, because the disjunctive guard is violated when either
conjunct fails and the weakest safe precondition requires both. A chained
dereference such as \texttt{a.b.c} contributes its definedness conditions ---
the non-nullness of each receiver in the chain --- which in WP terminology are
the $\mathrm{def}(e)$ side-conditions that must hold for the expression to
evaluate without fault. Division contributes a non-zero-divisor precondition,
and a division of the most-negative integer by minus one contributes an
overflow-avoidance precondition, because that single case overflows under
two's-complement semantics.

\subsection{What the Discipline Excludes}

The WP discipline is as important for what it excludes as for what it includes,
and the exclusions are where a naive pattern matcher would go wrong. A
condition that controls \emph{branching} rather than \emph{aborting} --- an
\texttt{if}/\texttt{else} in which both branches return a value, or the
condition of a ternary expression --- yields no WP-derived precondition on the
entry state, because neither branch aborts and so no entry condition is
required to avoid an abort. A pattern matcher that emitted such a condition as
a \texttt{requires} clause would misclassify internal control flow as a contract
requirement, over-constraining the method and producing a clause that a
verifier would reject. The engine therefore distinguishes aborting guards from
branching conditions and emits preconditions only for the former. This
distinction --- between a guard that says ``you may not call me unless $G$''
and a branch that says ``if $G$ then I do this, otherwise that'' --- is the
single most error-prone judgment in precondition inference, and getting it
right is the difference between a contract and a description of the
implementation.

\subsection{A Worked Example}

The distinction is best seen on a concrete pair of methods. Consider first a
defensive accessor:

\begin{lstlisting}
public static int firstOf(int[] arr) {
    if (arr == null) {
        throw new NullPointerException("arr");
    }
    if (arr.length == 0) {
        throw new IllegalArgumentException("empty");
    }
    return arr[0];
}
\end{lstlisting}

Both guards abort, so both contribute WP-derived preconditions. The engine
emits \texttt{requires arr != null} from the first guard and
\texttt{requires arr.length > 0} from the second (the negation of
\texttt{arr.length == 0}). The return statement contributes the postcondition
\texttt{ensures \textbackslash result == arr[0]}, and because the body reads but
never writes, the engine emits \texttt{assignable \textbackslash nothing}. The
inferred contract is exactly the contract a careful engineer would write, and
every clause discharges against the body: the array access \texttt{arr[0]} is
in bounds precisely because the two preconditions establish non-nullness and
non-emptiness.

Now consider a superficially similar method whose leading condition does not
abort:

\begin{lstlisting}
public static String sign(int x) {
    if (x < 0) {
        return "negative";
    } else {
        return "non-negative";
    }
}
\end{lstlisting}

Here the condition \texttt{x < 0} controls branching, not aborting: neither arm
throws, and both return a value. A pattern matcher that treated every leading
\texttt{if} as a guard would emit the unsound and nonsensical precondition
\texttt{requires x >= 0}, forbidding exactly the inputs the method is designed
to handle. The engine instead recognises that both arms return and emits no
precondition from this condition. It does, however, record the condition as a
\emph{path condition} for branch-sensitive postcondition inference, producing
the two guarded clauses
\texttt{ensures x < 0 ==> \textbackslash result.equals("negative")} and
\texttt{ensures x >= 0 ==> \textbackslash result.equals("non-negative")}. The
contrast between the two methods --- identical in surface shape, opposite in the
clauses they should yield --- is the reason the aborting-versus-branching
judgment is implemented as a distinct analysis step rather than folded into the
guard matcher.

\section{Postcondition Inference}
\label{sec:inf-post}

Postconditions are forward-flow predicates derived, under the SP discipline,
from return statements and field assignments. A return statement
\texttt{return e} contributes the clause \texttt{\textbackslash result == e},
the strongest-postcondition consequence relating the return value to the
returned expression. A field assignment \texttt{this.f = e} contributes the
clause \texttt{this.f == e}, with field references on the right-hand side
wrapped in \texttt{\textbackslash old(\dots)} to capture the SP convention that
the assigned value is computed from the pre-state. A method that constructs and
returns a new object contributes the postcondition
\texttt{\textbackslash result != null}, because object construction cannot
return null.

The treatment of mutable parameters follows the standard SP convention. Java
parameters can be reassigned within a method body, which would render a
postcondition referring to a parameter ambiguous: does \texttt{x} in the
postcondition mean the value passed in or the value after some internal
reassignment? SP resolves this by modelling each parameter as a fresh
entry-value variable and expressing postconditions in terms of that entry
value. In emitted JML the convention is rendered by using the parameter name
directly --- JML interprets a bare parameter name in a postcondition as its
entry value --- and \texttt{\textbackslash old(\dots)} for fields, which are
not frozen at entry the way parameters are.

\subsection{Branch-Sensitive Postconditions}

Branching receives finer-grained treatment than the textbook SP rule for
conditionals prescribes. The textbook rule combines the two branches of an
\texttt{if}/\texttt{else} by disjunction:
$\mathrm{SP}(S_1, P \wedge C) \vee \mathrm{SP}(S_2, P \wedge \neg C)$. Emitting
this disjunction directly would be correct but lossy: it discards the
information that links each return value to the condition that produces it. A
consumer of the specification --- a test generator trying to construct an input
that exercises a particular return path, or a verifier checking a particular
branch --- wants to know not merely that the result is one of two values but
which condition yields which. The engine therefore records, for each return
statement, the \emph{path condition} that reaches it, and emits independent
guarded clauses $C \Rightarrow \mathrm{post}_1$ and
$\neg C \Rightarrow \mathrm{post}_2$ rather than the collapsed disjunction.
Ternary expressions are decomposed identically, each arm becoming a guarded
clause. A lightweight symbolic executor maintains the path condition as it
walks the body, so that a return nested several conditionals deep carries the
conjunction of the enclosing guards.

Path conditions are simplified algebraically before emission, so that the
resulting JML is readable rather than a thicket of redundant conjuncts. The
simplifier applies the obvious arithmetic identities: $x \le 0 \wedge x < 0$
reduces to $x < 0$, and $x \le 0 \wedge x \ge 0$ reduces to $x = 0$. The
simplification is purely syntactic and conservative --- it never strengthens or
weakens a condition, only rewrites it to an equivalent shorter form --- so it
cannot introduce unsoundness, only improve legibility.

\subsection{The Symbolic Executor and Path Conditions}

The branch-sensitive postcondition inference of the previous subsection rests on
a lightweight symbolic executor, and its design is worth describing because it is
the component that makes guarded postconditions possible without incurring the
cost of full symbolic execution. A full symbolic executor explores every path
through a method, maintaining a symbolic state and a path condition for each, and
suffers path explosion on branch-heavy code and divergence on loops. The engine's
executor is deliberately partial: it tracks the path condition --- the conjunction
of branch guards under which a program point is reached --- but does not maintain a
full symbolic store, and it does not explore loop bodies symbolically, deferring
loops to the invariant heuristics.

For postcondition inference this partial tracking is exactly enough. When the
executor reaches a return statement, it has accumulated the conjunction of the
guards on the path to that return, and it emits the guarded clause pairing that
path condition with the returned expression. A return nested two conditionals deep
carries the conjunction of both enclosing guards; a return in the else-branch of a
conditional carries the negation of that conditional's guard. The executor handles
both \texttt{if}/\texttt{else} statements and ternary expressions by the same
mechanism, decomposing each into per-path returns rather than collapsing them into
a disjunction, which is the choice that preserves the link between each return
value and its enabling condition.

The cost of the partiality is that the executor cannot infer postconditions that
depend on the symbolic values of intermediate computations --- it tracks the path
condition but not a full model of how each variable's value evolves along the
path. A method whose return value depends on a sequence of intermediate
assignments, rather than directly on the path condition and the parameters, yields
a weaker postcondition than a full symbolic executor would. This is a deliberate
recall sacrifice for tractability, consistent with the engine's overall stance:
the partial executor handles the common case --- a return whose value and enabling
condition are both locally apparent --- cheaply and correctly, and defers the
hard case rather than paying the cost of full symbolic execution to capture it.
The path conditions the executor produces are simplified algebraically before
emission, as the previous subsection described, so that even a deeply nested return
yields a readable guarded clause rather than an unsimplified thicket of conjuncts.

\subsection{Soundness Gating by Return Type}

A class of plausible-looking postconditions is unsound under Java's
two's-complement integer arithmetic, and the engine gates these by return
type. The clause \texttt{\textbackslash result >= 0} is true for
\texttt{Math.abs(x)} and for \texttt{x * x} only under mathematical-integer
semantics; under two's-complement, \texttt{Math.abs(Integer.MIN\_VALUE)} is
negative and \texttt{x * x} overflows to a negative value for large \texttt{x}.
The engine emits the non-negativity postcondition only when the return type is
floating point, where the concern does not arise, and suppresses it for
\texttt{int} and \texttt{long}. This is a representative instance of a general
principle: a postcondition that is true in the mathematical idealisation but
false in the machine's actual arithmetic must be gated by the conditions under
which the idealisation holds, so that every emitted clause can withstand formal
verification under the operational integer semantics the validation layer
enforces.

\section{Loop Invariant Inference}
\label{sec:inf-loops}

Loops are the point at which WP/SP reasoning is incomplete by syntax alone. The
WP of a \texttt{while} loop with respect to a postcondition requires a loop
invariant, and no purely syntactic procedure can derive invariants in general,
since invariant inference is undecidable. The engine addresses the requirement
through four complementary heuristics, applied in turn until one succeeds.

The first is \emph{bound-based} invariant inference for counter loops. A loop
of the form \texttt{for (int i = a; i < b; i++)} maintains, throughout its
execution, the bound $a \le i \le b$, and the engine emits a direction-aware
bound: $i \ge a$ for an incrementing loop, $i \le a$ for a decrementing one,
together with the relationship to the upper bound. The direction-awareness
matters, because an invariant that is true for an incrementing loop is false
for a decrementing one, and a pattern matcher that ignored the step direction
would emit an invariant that fails at loop entry or exit.

The second is \emph{accumulator pattern detection}. A loop that maintains a
running sum, product, maximum, or minimum --- the canonical reduction patterns
--- admits an invariant relating the accumulator to the portion of the data
processed so far. A maximum-finding loop, for instance, maintains that the
accumulator is greater than or equal to every element seen so far and equal to
some element seen so far; the engine emits both the universal and the
existential clause, because the postcondition that the result is the maximum is
not provable from either alone. These quantified invariants are the engine's
most ambitious output and the ones that most stress the downstream verifier, a
point developed in Section~\ref{sec:inf-validation}.

The third is \emph{monotonicity inference} from update expressions. A variable
updated only by addition of a non-negative quantity is non-decreasing; a
variable updated only by a guarded decrement that stops at zero stays
non-negative. The engine walks the loop body's update statements and emits
monotonicity invariants where the update shape guarantees them.

Each of the first three heuristics corresponds to a recognisable family of loops,
and the correspondence is what keeps them principled rather than ad hoc. The
bound-based heuristic targets the counting loop, whose invariant is a relationship
between the counter and its bounds; it is the most common loop shape and the most
reliably handled. The accumulator heuristic targets the reduction loop, whose
invariant relates a running result to the processed prefix; it is the most
ambitious, producing the quantified invariants that the verifier extension exists
to discharge. The monotonicity heuristic targets loops whose interest is the
direction in which a variable moves rather than its exact value, and it produces
the weaker but still useful invariants --- non-decreasing, non-negative --- that
suffice when an exact relationship is out of reach. The three together cover the
loops whose invariants are determined by their counter, their accumulation, or
their update direction, which is the bulk of the loops in ordinary code.

The fourth, a fallback when the first three do not apply, is \emph{bounded
symbolic unrolling}: the loop is unrolled a small fixed number of times (three
in the current configuration) and the resulting straight-line code analysed for
an invariant that holds across the unrolled iterations. Unrolling is a fallback
rather than a primary heuristic because it is unsound as a source of invariants ---
a predicate that holds for three iterations need not hold for all --- so its output
is treated as a candidate to be checked rather than as a trusted invariant, and the
validation layer is what catches the candidates that fail to generalise. The
fallback exists because a weak candidate invariant that the verifier can check is
more useful than no invariant at all: it gives the discharge step something to
attempt, and a candidate that happens to generalise is recovered, while one that
does not is flagged. When all four
heuristics fail to produce a non-trivial invariant, the engine emits a
conservative weak invariant in preference to nothing, so that the verification
condition is at least partially discharged rather than left entirely open.
Across the loops in the evaluation corpus, the four heuristics together produce
a non-trivial invariant for roughly eighty-five percent of loops, applying
uniformly to \texttt{for}, \texttt{while}, and \texttt{for-each} constructs.

\subsection{Promotion of Invariants to Postconditions}

A quantified loop invariant that remains valid at loop termination can be
promoted to a method postcondition by substituting the loop counter with its
exit value. A summation loop whose invariant states that the accumulator equals
the sum of the first $i$ elements yields, on termination with $i$ equal to the
array length, the postcondition that the result equals the sum of the whole
array. This promotion is how established loop predicates flow into the
strongest postcondition of the surrounding method, and it is what allows the
engine to emit a meaningful postcondition for a method whose entire behaviour
is a loop. The promotion is sound only when the invariant genuinely holds at
exit, which the validation layer checks; a non-decreasing variant promoted as
if it were monotonically decreasing is exactly the kind of error the discharge
step catches.

\section{Frame Conditions and Purity}
\label{sec:inf-frame}

A specification is incomplete without a frame condition --- a statement of which
locations a method may modify --- because a postcondition that constrains the
return value says nothing about whether the method also scribbled on the heap. In
JML the frame condition is the \texttt{assignable} clause, listing the locations
the method is permitted to modify; \texttt{assignable \textbackslash nothing}
declares the method pure (modifying no heap location), and \texttt{assignable
\textbackslash everything} is the conservative default that constrains nothing.
The frame condition is essential to modular verification: when a verifier
replaces a called method with its contract, it must know which locations the call
may have changed, or it cannot preserve any fact about the unchanged ones.

The engine infers frame conditions by analysing the writes a method performs. A
method whose body contains no field assignment, no array store, and no call to a
method with a non-trivial frame is pure, and the engine emits \texttt{assignable
\textbackslash nothing}. A method that assigns to a field \texttt{this.f} has
\texttt{f} in its assignable set; a method that stores into an array parameter has
that array's elements in its set. The analysis is transitive through calls: a
method that calls another inherits the callee's frame, so a method that is pure
in its own body but calls a mutating method is not pure overall. The conservative
default for a call into an unanalysed method is that it may modify anything, which
is why an otherwise-pure method that calls an unsummarised library method receives
a conservative frame rather than \texttt{\textbackslash nothing}.

The empirical regularity that almost every method in the evaluation corpus is
pure --- the corpus being utility and value-type code --- has a consequence the
distribution work of Chapter~\ref{ch:distribution} exploits directly: since
\texttt{assignable \textbackslash nothing} is the overwhelming common case, the
distribution format makes it the default and omits it from the wire
representation, recovering it on read. The frame condition is thus a small clause
with an outsized role: it is necessary for sound modular verification, it is
trivial to infer for pure methods, and its ubiquity is the basis for a size
optimisation downstream. The discharge results of Chapter~\ref{ch:testgen}
reflect its tractability --- assignable clauses discharge at a higher rate than
any other clause kind, because purity is the easiest property to establish
against a body that visibly performs no writes.

\section{Interprocedural Propagation}
\label{sec:inf-inter}

Methods call other methods, and a specification that ignored calls would miss
the obligations they impose. The engine handles calls through procedure
summaries, the standard mechanism for modular WP/SP analysis. It proceeds in
two passes. The first analyses methods in dependency order --- callees before
callers, following a topological sort of the call graph --- and caches the
inferred specification of each method. The second pass propagates these cached
specifications across call sites: when inference encounters a call
$\mathtt{m}(a_1, \dots, a_k)$, it retrieves the cached pre- and postconditions
of \texttt{m} and instantiates them by formal-to-actual substitution, exactly
as the textbook WP rule for procedure call prescribes. A precondition that
\texttt{m} requires of its $j$-th parameter becomes, at the call site, a
precondition on the $j$-th actual argument.

Calls into the Java standard library are resolved against a curated
specification database covering common operations on \texttt{String},
\texttt{Math}, \texttt{List}, \texttt{Map}, and the standard utility classes.
The database is maintained by hand, which bounds its coverage but also bounds
the risk of propagating an incorrect library summary; a curated entry is
trusted, and a call into a library method with no entry is treated
conservatively. Calls into other unannotated methods --- application methods
with no cached specification and no database entry --- are treated as
uninterpreted: the callee may throw, may modify any accessible state, and
returns an unconstrained value. This conservative treatment is the limit case
of the procedure-summary mechanism, and it is the reason a small number of
methods in the evaluation corpus receive only an empty contract: they delegate
entirely to a library method for which no summary is available, and the engine
declines to invent one.

A small example fixes the mechanism. Suppose a method \texttt{checkedGet} calls a
helper \texttt{elementAt(list, index)} whose inferred contract requires
\texttt{index >= 0} and \texttt{index < list.size()}. When the engine analyses a
caller \texttt{int second(List<Integer> xs) \{ return elementAt(xs, 1); \}}, the
interprocedural analyser retrieves the helper's cached preconditions and
substitutes the actual arguments --- \texttt{xs} for the formal list and
\texttt{1} for the formal index --- to obtain \texttt{1 >= 0} (a tautology the
simplifier drops) and \texttt{1 < xs.size()}. The latter becomes a precondition of
\texttt{second}: the caller may invoke it only on a list of size at least two. The
obligation, invisible in \texttt{second}'s own body, is recovered from the helper's
contract by substitution, which is exactly the textbook weakest-precondition rule
for procedure call applied by the single pass. The example also shows the single
pass at its competent best --- an unconditional callee precondition lifted cleanly
to the caller --- which is the case the next chapter contrasts with the guarded and
polymorphic cases the single pass cannot handle.

The single-pass character of this propagation --- each callee's contract read
once and substituted at the call site without further transformation --- is
adequate for unconditional preconditions but blind to the branch-conditional
and polymorphic-dispatch obligations that Chapter~\ref{ch:compositional}
addresses. That chapter's contribution is precisely a measurement of what this
single pass leaves behind, and a second pass that recovers it. The present
chapter describes the single pass as the engine's baseline interprocedural
behaviour; the limitation is named here and quantified there.

\subsection{The Standard Library Specification Database}

The conservative treatment of unsummarised calls would cripple the engine on real
code, because real code calls the standard library constantly --- a method that
calls \texttt{Math.abs}, \texttt{String.length}, or \texttt{List.get} would, under
the uninterpreted-call policy, lose all information at the call site. The engine
mitigates this with a curated database of specifications for common standard-
library operations, covering \texttt{String}, \texttt{Math}, the collection
interfaces \texttt{List} and \texttt{Map}, and the boxed-primitive utility
classes. When the interprocedural analyser encounters a call into a library method
with a database entry, it uses the curated specification in place of the
conservative summary: a call to \texttt{Math.abs} on a floating-point argument
contributes the non-negativity of the result, a call to \texttt{String.length}
contributes a non-negative integer, a call to \texttt{List.get} contributes the
in-bounds precondition on its index.

The database is curated by hand rather than inferred, which is a deliberate
trade-off. Hand curation bounds the database's coverage --- only the entered
methods are summarised --- but it also bounds the risk: a curated entry is
trusted, having been written and reviewed by a human, so propagating it cannot
introduce the unsoundness an automatically inferred library summary might. The
database is, in effect, a small set of axioms the engine takes on faith, and its
size is the price of avoiding the alternative of inferring specifications for the
standard library (which would require its source, often unavailable) or treating
every library call conservatively (which would lose most cross-library
information). The three uninferable methods in the evaluation corpus of
Chapter~\ref{ch:testgen} are precisely those that delegate to a library method
with no database entry: the engine declines to invent a summary and emits an empty
contract, which is the correct conservative behaviour at the boundary of the
database's coverage.

\subsection{The Boundaries of Heuristic Inference}

It is as important to state what the engine does not infer as what it does, both
for honesty and because the boundaries explain the recall figures of
Chapter~\ref{ch:testgen}. The engine does not infer specifications that require
reasoning about the global heap shape --- properties relating objects reachable
through long reference chains, which need a shape analysis the pattern matcher does
not perform. It does not infer concurrency properties --- the absence of data
races, the correctness of a locking discipline --- because these are not local
syntactic properties and the engine reasons locally. It does not fully resolve the
interaction of inheritance and specification: a method that overrides another
should, under behavioural subtyping, honour the overridden method's contract, and
the engine's handling of inherited specifications is partial. And it does not infer
the complex algorithmic postconditions that characterise a method's deep semantic
purpose --- that a sort produces a sorted permutation, that a search returns the
correct index --- when those postconditions are not visible in the local structure
of returns and assignments.

These boundaries are the source of the recall gap: the postcondition recall of
78.1\% against documented intent (Chapter~\ref{ch:testgen}) is the complement of
the documented behaviours that lie outside what local pattern matching can reach.
The boundaries are not failures of the implementation but consequences of the
heuristic approach, and crossing some of them --- the algorithmic postconditions
in particular --- would require the dynamic observation or the symbolic execution
that the engine deliberately forgoes for scalability. Naming the boundaries
precisely is what allows the recall figure to be read correctly: not as a measure
of how often the engine is wrong, but as a measure of how much documented intent
lies beyond the reach of local structural inference.

\section{Formal Validation}
\label{sec:inf-validation}

Because the engine matches patterns rather than computing symbolic
transformers, individual clauses are not guaranteed sound by construction. The
WP/SP framing prescribes \emph{what} the engine looks for; it does not
guarantee that each emitted clause is provable against the actual code.
Soundness is recovered \emph{a posteriori} by discharging every emitted clause
through OpenJML's Extended Static Checker~\cite{openjml}. After inference, the
verifier may be invoked on the annotated source; each clause is discharged
independently, and clauses that fail verification are flagged for review rather
than silently retained. The validation pipeline is distributed as a
containerised toolchain to support reproducibility across operating systems,
and it underpins a regression suite of several hundred end-to-end verification
tests that gate every change to the inference engine against the formal oracle.

\subsection{Extending the Verifier}

The discharge layer uses a forked OpenJML build that emits SMT-LIB
\texttt{define-fun-rec} declarations for the \texttt{\textbackslash sum},
\texttt{\textbackslash product}, and \texttt{\textbackslash num\_of}
generalised quantifiers. Upstream OpenJML reports these constructs as ``not yet
supported,'' which would leave the loop-accumulator postconditions that the
engine works hardest to produce undischargeable --- recorded as unverifiable
rather than proved or refuted. The fork supplies recursive function definitions
that give the solver a handle on these quantifiers, converting a category of
``cannot check'' into ``checks'' or ``fails,'' and it is bundled with the
toolchain so that the validation pipeline remains reproducible. The extension
is a contribution in its own right: it is the difference between an inference
engine whose most ambitious output is unverifiable in principle and one whose
output is verifiable in practice.

\subsection{Strictness Configuration}

The strictness of the verifier is fixed at a configuration that was settled
empirically over the course of the project, because the choice materially
affects both what discharges and how long discharge takes. Four flags are set.
\texttt{-{}-code-math=safe} interprets the executable code's integer arithmetic
with operational two's-complement semantics and overflow checks, so that a
clause must hold under the arithmetic the machine actually performs.
\texttt{-{}-spec-math=bigint} interprets specification arithmetic with unbounded
mathematical integers, which is the right choice for the \emph{intent} a
specification expresses while \texttt{code-math=safe} guards the
\emph{implementation}. \texttt{-{}-arithmetic-failure=hard} treats an
arithmetic exception as a verification failure rather than a warning, so that
overflow is caught rather than tolerated. \texttt{-{}-nullable-by-default}
treats unannotated reference types as potentially null, the conservative
assumption. The combination of \texttt{code-math=safe} with
\texttt{spec-math=bigint} was found to be the configuration that maximises
genuine discharge without provoking solver timeouts; stricter combinations
(\texttt{spec-math=safe}, feasibility checking on every path) were measured to
cause the SMT solver to time out on the project's preamble, discharging fewer
clauses despite being nominally stronger. The default solver is Z3; alternative
versions are bundled and selectable, and an alternative solver (cvc5) is
available for cross-checking but is not used by default because it does not
consistently terminate on the project's preamble within the per-test timeout.

\subsection{The Two Layers Are Independent}

Manual validation and formal discharge are independent oracles, and the
distinction is important enough to state explicitly. Manual validation asks:
does the inferred clause agree with the method's documented intent? Formal
discharge asks: is the inferred clause consistent with the method's
implementation? A clause can pass either without the other. A clause may match
the Javadoc-documented intent perfectly yet fail discharge because of a subtle
overflow boundary the documentation does not mention; conversely, a clause may
discharge against the implementation yet be judged incomplete relative to the
documented intent, because it captures only part of what the method promises.
The two layers catch different errors --- manual validation catches
clauses that are formally consistent but miss the point, formal discharge
catches clauses that look right to a human but conflict with the code's actual
arithmetic or aliasing behaviour --- and running both reduces the risk that
high precision under one masks errors the other would surface. This independence
is exploited in the evaluation of Chapter~\ref{ch:testgen}, where the two are
reported as separate measurements rather than conflated into a single accuracy
figure.

\section{Class Invariants and History Constraints}
\label{sec:inf-invariants}

Beyond the per-method clauses, the engine infers class-level invariants ---
predicates that hold in every visible state of an object, before and after every
public method. A class invariant captures a property of the object's
representation that the class maintains as a whole, rather than a property of a
single method. The canonical case is a numeric field constrained to a range: a
class with a \texttt{size} field that is incremented and decremented but never
allowed to go negative has the invariant \texttt{size >= 0}, and the engine infers
it by checking that every assignment, increment, and decrement of the field
preserves the bound.

The inference is more delicate than it appears, because a field's invariant must
be preserved by \emph{every} method, and a single method that violates it refutes
it. The engine therefore examines all of a field's mutations across the class
before concluding an invariant holds. A decrement guarded by a positivity check
--- \texttt{if (size > 0) size{-}{-}} --- preserves \texttt{size >= 0}, and the
engine recognises the guard; an unguarded decrement does not, and the engine must
not emit the invariant, because a sequence of calls could drive the field
negative. An early version of the engine emitted the invariant on the strength of
the field's declared initial value alone, ignoring the decrements, and the
validation layer caught the resulting unsoundness --- a concrete instance of the
two-layer architecture catching an over-eager heuristic, and the reason the
class-invariant analyser now walks every mutation rather than trusting the
initialiser.

History constraints --- the \texttt{\textbackslash old} references in
postconditions --- are the per-method analogue of class invariants in their
subtlety. A method that increments a counter has the postcondition
\texttt{this.count == \textbackslash old(this.count) + 1}, relating the new state
to the old. The engine wraps a field reference on the right-hand side of an
assignment in \texttt{\textbackslash old} when the field is read before being
written in the same expression, capturing the strongest-postcondition convention
that the assigned value is computed from the pre-state. Getting the
\texttt{\textbackslash old} wrapping right is necessary for soundness: a
postcondition that omitted it would assert \texttt{this.count == this.count + 1},
a contradiction, and one that over-applied it would refer to pre-state values
where the body uses post-state ones. The wrapping is therefore applied precisely,
to field references whose value the postcondition draws from before the method's
writes, and the discharge results confirm the precision --- the
\texttt{\textbackslash old}-bearing postconditions discharge at a rate consistent
with the rest, indicating the wrapping is neither omitted nor overapplied.

\section{An End-to-End Example}
\label{sec:inf-example}

The individual analyses are best understood together, applied to a single method
that exercises preconditions, a loop with an accumulator, a promoted
postcondition, and the frame condition. Consider a method that sums the elements
of an array:

\begin{lstlisting}
public static long sum(int[] arr) {
    if (arr == null) {
        throw new NullPointerException("arr");
    }
    long total = 0;
    for (int i = 0; i < arr.length; i++) {
        total += arr[i];
    }
    return total;
}
\end{lstlisting}

The engine processes this method as follows. The precondition analyser
(Section~\ref{sec:inf-pre}) recognises the leading guard-and-throw idiom and
emits \texttt{requires arr != null}, the negation of the aborting guard. No
other guard is present, so no further precondition is emitted; in particular the
loop condition \texttt{i < arr.length} is a branching condition internal to the
loop, not an aborting guard, and contributes nothing to the precondition.

The loop-invariant analyser (Section~\ref{sec:inf-loops}) examines the
\texttt{for} loop. The counter \texttt{i} is initialised to zero and incremented
by one with an upper bound of \texttt{arr.length}, so the bound-based heuristic
emits the direction-aware invariant \texttt{loop\_invariant 0 <= i \&\& i <= arr.length}.
The body \texttt{total += arr[i]} matches the accumulator pattern: \texttt{total}
is a running sum over the processed prefix of the array, so the accumulator
heuristic emits the invariant
\texttt{loop\_invariant total == (\textbackslash sum int k; 0 <= k \&\& k < i; arr[k])},
relating the accumulator to the sum of the first \texttt{i} elements. This is the
quantified invariant whose discharge motivates the \texttt{define-fun-rec}
extension to OpenJML described in Section~\ref{sec:inf-validation}; without that
extension the verifier would record it as unsupported.

The invariant-promotion step (Section~\ref{sec:inf-loops}) observes that the sum
invariant remains valid at loop termination, where \texttt{i} equals
\texttt{arr.length}, and substitutes the exit value of the counter to obtain the
method postcondition
\texttt{ensures \textbackslash result == (\textbackslash sum int k; 0 <= k \&\& k < arr.length; arr[k])}.
This is the strongest postcondition the engine can establish for the method, and
it is exactly the specification a human would write: the result is the sum of the
whole array. The frame analyser observes that the body writes only to the local
\texttt{total} and \texttt{i}, never to a heap location, and emits
\texttt{assignable \textbackslash nothing}.

The assembled contract is therefore:

\begin{lstlisting}
//@ requires arr != null;
//@ ensures \result ==
//@   (\sum int k; 0 <= k && k < arr.length; arr[k]);
//@ assignable \nothing;
public static long sum(int[] arr) { ... }
\end{lstlisting}

Every clause discharges. The precondition guards the array dereference inside the
loop; the loop invariants are inductive (preserved by the body and established on
entry) and combine with the exit condition to prove the postcondition; and the
frame condition holds because the body touches no heap location. The example
illustrates the division of labour the chapter has described: precondition
inference from the guard, loop-invariant inference from the counter and
accumulator, postcondition inference by promotion of the invariant, and frame
inference from the absence of writes --- each a separate analysis, justified by a
distinct correspondence to the WP/SP calculi, and combining into a contract that
the validation layer certifies. It also illustrates the return type's role in
soundness: the accumulator is declared \texttt{long}, and had it been declared
\texttt{int} the engine would have to weigh the overflow risk of summing a large
array, gating any non-negativity claim by the considerations of
Section~\ref{sec:inf-post}.

\section{Method Categorisation}
\label{sec:inf-cat}

To enable systematic evaluation across diverse implementations, the engine
assigns each method to one of six categories, extending the method-stereotype
taxonomy of Dragan et al.~\cite{dragan2006reverse}. The categories are
\emph{accessors and mutators} (methods that read or write a single field),
\emph{control flow} (methods organised around conditionals and loops),
\emph{factory and delegate} (methods that construct objects or forward to other
methods), \emph{utility} (static helper methods), \emph{state modification}
(methods that mutate object state under conditional logic), and \emph{other}
(event handlers and callbacks). When more than one category's pattern matches a
method, the most specific is chosen. The categorisation is not used by the
inference engine itself --- inference is category-blind --- but it structures
the evaluation, allowing the studies that follow to report where inferred
specifications are most accurate and where they most improve downstream tasks.
The taxonomy's inter-rater reliability, established by two independent raters
categorising the evaluation corpus, is high (Cohen's $\kappa = 0.87$), which
gives confidence that the per-category results of later chapters rest on a
stable partition rather than on idiosyncratic labelling.

The categorisation also exposes a structural feature of the evaluation corpus
that recurs in the threats-to-validity discussions of later chapters: the
\emph{other} category --- event handlers, I/O operations, callbacks --- is
empty in the Apache Commons Lang corpus, because that library is a collection
of pure utility and value types and contains no methods of that kind.
Experience on broader Java code suggests this category is the hardest to
specify heuristically, since its methods often have effects that escape the
local syntactic patterns the engine recognises, and its absence from the
controlled corpus is a limitation on the external validity of the results that
the relevant chapters name explicitly.

\subsection{How Inference Behaviour Varies by Category}

Although inference is category-blind, the engine's behaviour varies
systematically across the categories, and understanding the variation explains the
per-category accuracy results of Chapter~\ref{ch:testgen}. For accessors and
mutators, inference is at its most reliable: a getter's postcondition is the field
it returns, a setter's is the assignment it performs, and both are captured by the
most direct postcondition patterns, which is why this category attains the highest
precision. For control-flow methods, precondition inference dominates: their guards
yield \texttt{requires} clauses reliably, and their branch structure yields guarded
postconditions, so the engine produces strong specifications whenever the branches
return values the patterns recognise. For utility methods, the picture is mixed:
their input-validation guards yield good preconditions, but their computational
postconditions --- what the utility actually computes --- often lie beyond the
local patterns, which is why their postcondition recall trails their precondition
recall.

For state-modification methods, inference is at its hardest, and the category's
lower specification strength reflects it. A method that mutates object state under
conditional logic has a frame condition and a set of guarded state-change
postconditions that the engine must assemble from the assignments scattered through
its branches, and capturing the complete effect --- every field that might change,
under every condition --- is the inference task most prone to incompleteness. The
engine captures the assignments it can localise and emits guarded postconditions
for them, but a method whose state change is distributed across many branches or
mediated by helper calls yields a partial rather than a complete specification. For
factory and delegate methods, inference depends heavily on interprocedural
propagation: a factory's postcondition is that it returns a non-null constructed
object, captured directly, but a delegate's behaviour is its callee's behaviour,
captured only if the callee has a cached or curated specification --- which is why
the three uninferable methods of the corpus are delegates to an unsummarised
library method. The category variation in inference behaviour is thus the upstream
cause of the category variation in both specification accuracy
(Chapter~\ref{ch:testgen}) and test-generation benefit: the categories the engine
specifies most completely are not always the categories where specifications help
most, because the benefit depends also on the baseline, but the categories the
engine specifies least completely --- state modification, complex delegation --- are
consistently the ones where both accuracy and benefit are bounded by what local
inference can reach.

\section{Operation, Logging, and Diagnostics}
\label{sec:inf-operation}

A research instrument that is run thousands of times over evolving code needs to be
observable, and the engine's logging and diagnostics, while not an analytical
contribution, are part of what made the empirical studies possible. The engine logs
at two levels: a progress level reporting which files it processes and how many
specifications it produces, suitable for a console during a run, and a detailed
level recording, for each method, which patterns fired, which clauses were emitted,
and which were suppressed by the soundness gates. The detailed log is what allows a
maintainer to understand why a particular method received the specification it did
--- why a clause was emitted, or, more often usefully, why an expected clause was
not --- and it is the first recourse when a regression test reveals that a change to
one analyser altered another's output.

The diagnostic value of the logging is most apparent at the interprocedural and
compositional boundaries, where a method's specification depends on others' and a
surprising result is most likely to have a non-local cause. When a call site fails
to propagate an expected precondition, the log records whether the callee was
resolved, whether its specification was cached, whether the substitution succeeded,
and whether the result was already present --- the same diagnostic breadcrumbs that
the compositional measurement of Chapter~\ref{ch:compositional} aggregates into its
call-site resolution statistics. The instrumentation that supports day-to-day
maintenance is thus the same instrumentation that supports the measurement: a count
of how often each step succeeds or fails per call site is both a debugging aid and a
datum, and the compositional study's report of arity-match rates is, in effect, the
detailed log summarised over a corpus.

This observability reflects the engine's status as evolving research software. A
heuristic engine whose pattern repertoire grows continually cannot be understood by
reading its code alone, because its behaviour on a given method is the combined
effect of many patterns interacting; the log is how that combined effect is made
legible. The investment in diagnostics is, in this sense, of a piece with the
two-tiered testing and the file-size discipline: all three are the
software-engineering practices that let a heuristic engine remain understandable and
correct as it accumulates the patterns that broaden its reach.

\section{Summary}

JML-Inferrer infers JML preconditions, postconditions, loop invariants, and
frame conditions from Java source by pattern matching organised around the
weakest-precondition and strongest-postcondition calculi. It avoids symbolic
WP/SP computation, which does not scale, and uses the calculi instead as a
design discipline that fixes its pattern repertoire and clause-generation
rules. It propagates specifications across method calls by single-pass
procedure summarisation. It recovers the soundness that pattern matching does
not guarantee by discharging every emitted clause through a forked, quantifier-
extended OpenJML, treating inference and verification as two independent layers
in which the first proposes and the second disposes. The chapters that follow
put this output to work: Chapter~\ref{ch:testgen} measures its value to a
language-model test generator and reports its accuracy under both the manual
and the formal oracle; Chapter~\ref{ch:distribution} carries it across the
source-to-binary boundary; and Chapter~\ref{ch:compositional} measures and
recovers the compositional information that the single-pass propagation of
Section~\ref{sec:inf-inter} leaves behind.
